\section{Computational Pipeline}

This appendix describes the computational procedures used to verify
the structural claims in the paper.

\subsection{Graph construction}

The graph $\Gsixty$ is constructed from the $120$ flags of the
dodecahedron by identifying opposite configurations to obtain $60$
chamber states. Adjacency is generated by local transport moves.

The resulting graph is verified to be:
\begin{itemize}
\item connected,
\item $4$-regular,
\item with $60$ vertices and $120$ edges.
\end{itemize}

\subsection{Quotient computation}

Successive quotients
\[
\Gsixty \to \Gthirty \to \Gfifteen
\]
are computed by identifying vertices under involutive symmetries.

The graph $\Gfifteen$ is verified to be isomorphic to
$L(\mathrm{Petersen})$ by comparison of adjacency structure.

\subsection{Signed lift and cocycle}

A signing $\sigfun : E(\Gfifteen) \to \{\pm1\}$ is computed from the
transport rules. The induced double cover is verified to reproduce
$\Gthirty$.

Switching equivalence is implemented via vertex sign assignments, and
the minimal support is obtained by exhaustive search over the switching
orbit.

\subsection{Minimal support search}

All $2^{14}$ switching configurations (fixing one vertex) are evaluated.
For each configuration, the number of negative edges is recorded.

The minimum support size is found to be $6$, and all minimal supports
are enumerated.

\subsection{Sector construction}

The sector--edge incidence matrix $M$ is constructed from the transport
data. Matrix products are computed explicitly and verified to satisfy
\[
MM^{\mathsf T} = A^3 + 2A^2 + 2I,
\]
where $A$ is the adjacency matrix of $\Gfifteen$.

\subsection{Implementation}

All computations are carried out using Python scripts operating on
explicit adjacency and incidence data. The computations are entirely
deterministic and can be reproduced from the data and procedures
described above.
